% !TEX program = xelatex
% Raspberry Pi Xinu マルチコア化レポート
%
% ボードを追加するには:
%   1. boards/ に rpi4.tex / rpi5.tex を用意する（boards/rpi3.tex が雛形）
%   2. 下の「ボード別」節の \input を有効にする
%   3. \boardrow{...} を横断比較表に足す（各ボード .tex の末尾で定義済みの
%      \rpiXXXsummary を呼ぶだけで済むようにしてある）
% 本文の共通部（設計・測定方法・結論）はボードに依存しない。

\documentclass[11pt,a4paper]{article}

\usepackage[a4paper,margin=22mm]{geometry}
\usepackage{xeCJK}
\usepackage{fontspec}
\usepackage{amsmath,amssymb}
\usepackage{listings}
\usepackage{xcolor}
\usepackage{booktabs}
\usepackage{tabularx}
\usepackage{longtable}
\usepackage{array}
\usepackage{hyperref}
\usepackage{enumitem}
\usepackage{makecell}
\usepackage{graphicx}

% 測定待ちの箇所を目立つ形で残すマクロ（数値が出たら本文を差し替える）。
\newcommand{\TODO}[1]{\par\medskip\noindent%
  \colorbox{yellow!25}{\parbox{0.97\linewidth}{%
  \textbf{［測定待ち］}\ #1}}\par\medskip}

\setCJKmainfont{Hiragino Mincho ProN}
\setCJKsansfont{Hiragino Sans}
\setCJKmonofont{Hiragino Sans}
\setmonofont{Menlo}[Scale=0.82]

\hypersetup{
  colorlinks=true,
  linkcolor=black,
  urlcolor=blue!60!black,
  citecolor=blue!50!black,
  pdftitle={Raspberry Pi 上の Embedded Xinu のマルチコア化},
  pdfauthor={Yasushi Kodama (児玉靖司)}
}

% ---- listings ----------------------------------------------------------
\definecolor{codebg}{RGB}{248,248,242}
\definecolor{codecmt}{RGB}{120,120,120}
\definecolor{codekw}{RGB}{30,80,160}
\lstset{
  basicstyle=\ttfamily\scriptsize,
  backgroundcolor=\color{codebg},
  commentstyle=\color{codecmt}\itshape,
  keywordstyle=\color{codekw}\bfseries,
  breaklines=true,
  frame=single,
  framesep=3pt,
  rulecolor=\color{black!25},
  showstringspaces=false,
  captionpos=b,
  language=C,
}

% ---- ボード横断で使う共通マクロ ----------------------------------------
% 各ボード節はこの体裁に従う。新ボードもこれを使えば見た目が揃う。
\newcommand{\bench}[1]{\texttt{#1}}
\newcommand{\route}[1]{\texttt{#1}}
% 速度向上の書式（測定値は「比」で報告する。理由は §3.3）
\newcommand{\spd}[1]{$\times$#1}

\title{\textbf{Raspberry Pi 上の Embedded Xinu のマルチコア化}\\[2mm]
       \large ワーカープール型 SMP の設計と実測\\[1mm]
       \normalsize 何が速くなり、何が速くならないか}
\author{児玉靖司\\\small 法政大学経営学部\\\small \texttt{yass@hosei.ac.jp}}
\date{2026年7月16日}

\begin{document}
\maketitle

\begin{abstract}
\noindent
Raspberry Pi 上でベアメタル動作する Embedded Xinu を、搭載された 4 つの
CPU コアを使うように拡張した。設計は、コア 0 が既存の単一コア OS
（スケジューラ・アクター・ネットワーク・ウィンドウシステム）を一切変更せずに
実行し続け、残るコアを純計算専用のワーカーとして起こす「ワーカープール型」
である。対称型 SMP を採らないのは、既存ランタイムが非リエントラントであり、
共有 ready キューが実績のある機能を破壊するためである。

本レポートの主眼は、速くなったという報告ではなく、\textbf{何が速くなり、
何が速くならないかを実測で切り分けた}点にある。計算律速の処理は
最大 \spd{3.99}（ほぼ理想値）まで向上した一方、メモリ書込律速の処理は
\spd{0.95}--\spd{1.35} に留まり、並列化の意味がないことが分かった。
この結果は「ウィンドウシステムを 4 コアで速くできるか」という問いに
否と答える根拠となった。またワーカープールの起動コストは実測 7--11\,\textmu s
であり、これを下回る仕事は並列化すると\textbf{かえって遅くなる}
（実測 \spd{0.38}）。この 2 つの境界が、並列化すべき対象を決める。

さらに、この「メモリ律速はスケールしない」という結論が\textbf{D-cache OFF に
固有}であることを、D-cache を有効化する実験で実測した（第 \ref{sec:dcache} 節）。
D-cache ON では、作業セットがキャッシュに収まる限りメモリ書込は理想的に
スケールし（64\,KB で \spd{4.01}）、超えると帯域律速へ回帰する（1\,MB で
\spd{1.40}）。ただし実フレームバッファはキャッシュ不可であり、かつ全画面
描画はキャッシュを超えるため、ウィンドウシステムの結論自体は覆らない。
D-cache の有効化は Cortex-A53 で set/way invalidate・SMPEN・DMA コヒーレンシを
要し、とりわけ USB ドライバがスタック上のバッファを DMA する構造が
最大の障害であることを明らかにした。

Raspberry Pi 3 の実装と測定を第 \ref{sec:rpi3} 節、D-cache 実験を
第 \ref{sec:dcache} 節に示す。さらに Raspberry Pi 4（Cortex-A72）でも同一の
ハーネスで測定し、「計算はスケールし、メモリはスケールしない」という結論が
微アーキテクチャを超えて一般化することを確認した（第 \ref{sec:cross} 節）。
\end{abstract}

\tableofcontents
\newpage

% =======================================================================
\section{はじめに}

Raspberry Pi に搭載される BCM2837 / BCM2711 / BCM2712 は、いずれも 4 コアの
ARM プロセッサを持つ。しかし Embedded Xinu は単一コア OS であり、
これらのボード上では 1 コアしか使わず、残る 3 コアは起動されないまま
放置されていた。本作業の目的は、この遊休コアを実際の計算に使うことである。

本レポートは次の順で述べる。第 \ref{sec:design} 節で全ボード共通の設計方針を、
第 \ref{sec:method} 節で測定方法を述べる。第 \ref{sec:rpi3} 節以降が
ボード別の実装と実測である。第 \ref{sec:cross} 節で横断比較を行い、
第 \ref{sec:conclusion} 節で結論を述べる。

% =======================================================================
\section{背景と関連研究}
\label{sec:related}
% ======================================================================
% related_work.tex — 背景と関連研究
% ======================================================================

Xinu は Douglas Comer が教科書『Operating System Design: The Xinu Approach』の
ために設計した、単一プロセッサ前提の軽量教育用 OS である。その ARM / MIPS 系
組込み移植 \textbf{Embedded Xinu} は Marquette 大学の Dennis Brylow のグループが
主導し、複数大学の OS 教育で使われている。本作業もこの Embedded Xinu の
Raspberry Pi 3 移植を出発点とする。

\subsection{既存研究：XinuPi3（Marquette、教育目的の対称型 SMP）}

Raspberry Pi 3 の 4 コア化そのものは、Xinu の本家グループが既に発表している。
\textbf{XinuPi3}\footnote{P. Bansal, R. Latinovich, T. Lazar, P. J. McGee,
D. Brylow, ``XinuPi3: Teaching Multicore Concepts Using Embedded Xinu,''
CSERC '17（6th Computer Science Education Research Conference）, Helsinki, 2017.
\url{https://epublications.marquette.edu/mscs_fac/634/}} は Embedded Xinu を
Pi 3 へ移植し、\textbf{対称型マルチプロセッシング（full SMP）}を実現した。
その主眼は\textbf{教育}にあり——「マルチコア概念を、真に並行なハードウェアで
初学者に教える」ため——コアの起動、\textbf{キャッシュコヒーレンシ}、
\textbf{コア間相互排他（ロック）}、マルチコアスケジューリングを扱う。
論文自身が「ベアメタルの教育用 OS でマルチコアを支えるものは他に無い」と述べ、
その\textbf{存在}と\textbf{教材としての価値}を貢献としている。同グループは
続けて、Pi 3 を用いた OS 教育・アセンブリでの並列計算教育の実践も報告して
いる\footnote{P. J. McGee, R. Latinovich, D. Brylow, ``Using Embedded Xinu
and the Raspberry Pi 3 to Teach Operating Systems,'' IEEE IPDPS Workshops
(EduPar), 2020.}。

\subsection{本研究との違い}

本レポートの研究は、XinuPi3 と\textbf{目的・方式・貢献のいずれも異なる}。

\begin{table}[htbp]
\centering
\caption{XinuPi3（既存）と本研究の違い}
\small
\begin{tabularx}{\textwidth}{@{}lXX@{}}
\toprule
 & \textbf{XinuPi3（Marquette）} & \textbf{本研究} \\
\midrule
目的 & 教育——マルチコア概念を教える教材 & 測定駆動の工学——稼働中の機能豊富な OS の並列性能を、退行させずに最大化 \\
SMP 方式 & \textbf{対称型（full SMP）} & \textbf{非対称ワーカープール型}。コア 0 は既存 OS を無改変で走らせ、コア 1--3 を計算ワーカーに \\
キャッシュ & D-cache 有効化＋コヒーレンシ処理 & D-cache は\textbf{意図的に OFF}（volatile+dsb でロックフリー）。その上で「ON にすると何が変わるか」を独立実験（第\ref{sec:dcache}節） \\
同期 & コア間ロック（相互排他） & \textbf{ロックフリー}。full SMP を採らないので cross-core ロックが不要 \\
対象カーネル & 教育用の素の Xinu & WiFi ad-hoc メッシュ・WM・AIPL アクター・オンデバイス JIT が\textbf{稼働中で不変条件} \\
貢献 & マルチコア対応の\textbf{存在}と教材価値 & \textbf{何が速くなり何が速くならないか}の定量的切り分け＋設計方法論 \\
\bottomrule
\end{tabularx}
\end{table}

\noindent
最も本質的な違いは、\textbf{対称型 SMP か、非対称ワーカープールか}である。
XinuPi3 は full SMP を教育目標として\textbf{採る}。本研究は、既存ランタイム
（アクター・ネット・WM）が非リエントラントであるという現実から出発し、
full SMP は実績ある機能を全面的に作り直すリスクを負うと判断して\textbf{採らない}。
そしてこの判断が正しいことを、二重に示した——実機実測で worker-pool が計算律速で
\spd{3.99}（ほぼ理想値）を出すこと（第\ref{sec:rpi3}節）と、設計空間 GA が
worker\_pool を full\_smp に対する champion として選ぶこと（第\ref{sec:ga}節）で
ある。XinuPi3 が「4 コアを使えること」を示したのに対し、本研究は
「4 コアで\textbf{何を}どこまで使うべきか、そして何は使っても無駄か」を測る。

第二の違いは方法論にある。本研究は、(a) 実機での定量測定（計算 vs メモリ律速、
プール往復の損益分岐、キャッシュ構成の支配性）、(b) D-cache 有効化という制御された
比較実験、(c) AI レビュアーによる設計空間 GA と実測の照合——という、\textbf{測定と
探索を軸にした}接近を採る。教育用途を主眼とする既存研究とは、問いの立て方が異なる。


% =======================================================================
\section{共通設計：ワーカープール型 SMP}
\label{sec:design}

\subsection{なぜ対称型 SMP を採らないか}

素朴な発想は、ready キューを全コアで共有する対称型 SMP である。しかし
既存の Xinu ランタイム（AIPL アクター、ネットワークスタック、USB、
ウィンドウマネージャ）は\textbf{非リエントラント}であり、複数コアから
同時に呼ばれることを想定していない。共有 ready キューを導入すれば、
実機で動作実績のあるこれらの機能を全面的に作り直す必要が生じ、
リスクが利得を上回る。

そこで本設計では次の非対称構成を採る。

\begin{itemize}[leftmargin=1.4em]
  \item \textbf{コア 0} は既存の単一コア OS をそのまま実行する。
        ready キューはコア 0 の専有であり、スケジューラ・アクター・
        ネットワーク・描画は\textbf{一切変更しない}。
  \item \textbf{コア 1--3} は専用の計算ワーカーとして起動する。低消費電力の
        WFE で待機し、投入されたジョブ（仕事の範囲）を実行して完了を通知する。
        OS には触れず、割り込みも受けず（IRQ/FIQ マスク）、メモリ確保も
        行わないため、カーネルと競合し得ない。
\end{itemize}

この構成は、既存機能への退行リスクが構造的にゼロであるという点で優れている。
コア 0 の挙動が変わらないからである。

\subsection{ロックフリー調停が成立する理由}

ワーカーとコア 0 の調停は、単一のジョブメールボックス（\texttt{volatile} な
配列群）と \texttt{dsb} バリアのみで行い、スピンロックも
LDREX/STREX も使わない。これが成立するのは、この移植が\textbf{D-cache を
OFF にしている}ためである。すなわち全てのデータアクセスが RAM に直行するので、
コア間のキャッシュコヒーレンシ問題が原理的に生じない。

これは設計上の仮定であり、実機で検証した（第 \ref{sec:rpi3} 節）。
なお D-cache を有効化すると、この仮定は崩れ、ジョブメールボックスには
実際のキャッシュ管理が必要になる。

\subsection{ジョブ引数は呼び出し側スタックに置く}
\label{sec:ud}

ジョブの引数をファイルスコープの変数で渡してはならない。コア 0 は
プリエンプティブであり、かつプールの利用者は複数スレッド（HTTP サーバ
スレッドと AIPL アクタースレッド）に及ぶ。ゆえに次の競合が起こり得る。

\begin{enumerate}[leftmargin=1.6em]
  \item スレッド A が引数を書く。
  \item A がプリエンプトされる。
  \item スレッド B が引数を上書きし、プールを取って実行する。
  \item A が再開し、\textbf{B の引数で計算して誤った答えを返す}。
\end{enumerate}

本設計では API に不透明ポインタ \texttt{ud} を設け、引数を呼び出し側の
スタックに置くことで、この競合を構造的に起こり得なくした。ロックで
塞ぐより安全であり、副次的に性能も向上した（第 \ref{sec:rpi3-volatile} 項）。

\subsection{プールの排他とフォールバック}

ジョブメールボックスは 1 組しかないため、プールは同時に 1 人しか使えない。
既に使用中の場合、後続の呼び出し元は\textbf{自分の全チャンクをその場で実行}する。
待たせないので デッドロックせず、並列化の利得を諦めるだけで答えは常に正しい。

同様に、ワーカーが応答しない場合はコア 0 がそのチャンクを引き取る。
したがって\textbf{ワーカーの不調が OS を停止させることはない}。

% =======================================================================
\section{測定方法}
\label{sec:method}

\subsection{時計：割り込みではなくフリーランニングカウンタを読む}
\label{sec:clock}

当初、経過時間の測定には Xinu の \texttt{clktime}/\texttt{clkticks} を用いていたが、
\textbf{どの測定も 0\,ms を返した}。原因はこれらがタイマ\textbf{割り込み}で
加算される変数であり、その割り込みが安定して発火していないことにある。
これはマルチコア化以前のカーネルでも同様であり、本作業が持ち込んだ
問題ではない。

そこで System Timer の 1\,MHz フリーランニングカウンタを直接読む
\texttt{clkcount()} に切り替えた。割り込みに依存しないため停止せず、
分解能も 1\,ms から 1\,\textmu s へ 1000 倍向上した。

\subsection{正しさの同時検証}

全てのベンチマークは、同一の仕事を\textbf{逐次で 1 回、並列で 1 回}実行し、
両者の答えが一致するかを \texttt{agree} として毎回報告する。速いが誤った
結果を、速いという理由で見逃さないためである。本レポートに載せた全測定で
\texttt{agree = yes} である。

さらに、分割ロジック（剰余類による仕事の配分）は実機に投入する前に、
カーネルのソースから当該関数をそのまま抜き出してホスト上で総当たり検証した。
この検証により、実機に載せる前に実際に 2 件の誤りを発見・修正している
（第 \ref{sec:rpi3-bugs} 項）。

\subsection{絶対時間ではなく比を報告する}
\label{sec:ratio}

測定中に、\textbf{CPU クロックが一定でない}ことを発見した。再起動直後の
最初の測定だけが速く、以降は安定した遅い値に落ち着く。その比は実測 2.33 で、
Raspberry Pi 3 B+ の $1.4\,\mathrm{GHz} \div 600\,\mathrm{MHz} = 2.33$ と
一致する。ベアメタル Xinu は DVFS を制御しないため、ファームウェアが
クロックを変更しているものと考えられる。

したがって\textbf{絶対時間は再起動からの経過に依存し、比較に使えない}。
一方 \textbf{速度向上比は安定}している。1 コア測定と 4 コア測定を同一
リクエスト内で連続実行しているため、クロックの影響が相殺されるからである。
実際、繰り返し測定で比は完全に一致した（例：\spd{3.99} が 3 回とも同値）。
本レポートは比を主たる指標とし、絶対時間は状態を明記した上で参考値として示す。

% =======================================================================
\section{Raspberry Pi 3 (BCM2837)}
\label{sec:rpi3}
% ======================================================================
% boards/rpi3.tex — Raspberry Pi 3 (BCM2837) のマルチコア化
%
% ★ このファイルは rpi4.tex / rpi5.tex の雛形でもある。
%    節構成（実装 → 起動 → 適用 → 実測 → 発見したバグ）を揃えると
%    main.tex の横断比較表がそのまま使える。
%    \section{} は main.tex 側に置いてあるので、ここは \subsection から始める。
% ======================================================================

\subsection{対象と前提}

対象は Raspberry Pi 3 B+（BCM2837、Cortex-A53 $\times$ 4、1.4\,GHz）上の
Embedded Xinu（リポジトリ \texttt{xinu-raz/xinu}）である。このボードは
WiFi・有線 LAN・HDMI ウィンドウマネージャ・AIPL アクターランタイム・
JIT が実機で動作している一方、\texttt{include/rcu.h} に
「\textit{The arm-rpi3 port parks the secondary cores (MPIDR guard)}」と
あるとおり、\textbf{4 コアのうち 3 コアを起動せずに放置}していた。

本移植の前提として重要なのは、このポートの MMU 設定である。
\texttt{system/platforms/arm-rpi3/mmu.c} は identity map（VA = PA）を張り、
\textbf{MMU と I-cache は有効、D-cache は無効}のままにしている。
実機で読み出した \texttt{SCTLR} がこれを裏付ける。

\begin{lstlisting}[language=,caption={実機 \route{/api/mmu} の応答（D-cache OFF の確認）}]
enabled=1
sctlr=0x00c51839 M=1 C=0 I=1 Z=1
ttbr0=0x0011c000
\end{lstlisting}

\noindent
$C = 0$、すなわち D-cache は OFF である。よって第\ref{sec:design}節の
ロックフリー調停の前提が、この実機で成立していることが確認できた。

\subsection{コア起動：Pi 3 固有の差分}
\label{sec:rpi3-bringup}

Pi 4 / Pi 5 は AArch64 で PSCI やスピンテーブルによりコアを起動するが、
Pi 3 は 32bit（\texttt{arm\_64bit=0}）で動作しており、機構が異なる。
32bit のファームウェアは\textbf{コア 0 しか起動せず}、コア 1--3 は armstub7 の
中で自分専用のメールボックス 3 を監視して待機している。そこへ入口アドレスを
書き込み、\texttt{sev} で起こす。

\begin{lstlisting}[caption={ARM ローカルのコア別メールボックス 3（\texttt{system/platforms/arm-rpi3/smp.c}）}]
/* ARM-local per-core mailbox 3 SET register.  The 32-bit Pi 3 firmware leaves
 * cores 1-3 spinning in armstub7 on exactly this word; writing a jump target
 * and issuing SEV releases the core to that address.  This MMIO page is
 * already Device-mapped by mmu.c (0x40000000..0x40FFFFFF). */
#define ARM_LOCAL_MBOX3_SET(core) \
    ((volatile unsigned long *)(0x4000008CUL + 0x10UL * (unsigned long)(core)))
\end{lstlisting}

\noindent
この MMIO 領域は既に mmu.c が Device 属性で写像済みであったため、追加の
写像は不要であった。起こされたコアは \texttt{start.S} に追加した
\texttt{\_smp\_start} トランポリンに入り、SVC モードへ正規化し、
IRQ/FIQ をマスクし、コア専用スタックを設定して C の入口へ入る。そこで
コア 0 と\textbf{同一の}MMU 設定（コア 0 が構築した L1 テーブルを共有し、
MMU と I-cache を有効化、D-cache は無効のまま）を適用する。設定を揃えるのは
体裁の問題ではなく、4 コアの時間比較を公平にするためである。

なお、コア解放の手掛かりを 2 系統用意した。armstub7 のメールボックスと、
\texttt{start.S} 側の \texttt{smp\_release[]} である。後者はファームウェアが
4 コアとも \texttt{\_start} に流し込む実装だった場合の保険である。この配列は
\textbf{意図的に \texttt{.bss} ではなく \texttt{.data} に置いた}。コア 0 が
\texttt{.bss} をゼロクリアしている最中に、待機中のコアが古い値を読むことを
防ぐためである。

\begin{table}[htbp]
\centering
\caption{起動が正しく構成されたことの確認（リンク結果）}
\small
\begin{tabular}{@{}lll@{}}
\toprule
シンボル & アドレス & セクション \\
\midrule
\texttt{\_smp\_start}   & \texttt{0x00008280} & \texttt{.text} \\
\texttt{smp\_release}   & \texttt{0x0010e980} & \texttt{.data}（\texttt{D}）\\
\texttt{smp\_stacktop}  & \texttt{0x0010e990} & \texttt{.data}（\texttt{D}）\\
\bottomrule
\end{tabular}
\end{table}

\subsubsection*{実機での結果}

最大の不確実性は「32bit ファームウェアが本当にこのメールボックスでコアを
解放するか」であった。実機で\textbf{確認された}。

\begin{lstlisting}[language=,caption={\route{/bench} の応答（4 コアが起動している）}]
SMP bench kind=nqueens
cores_online = 4
\end{lstlisting}

\noindent
仮に外れていても、起動待ちは短くバウンドしてあり、応答しないコアは
offline 扱いで 1 コアにフォールバックして起動は継続する設計とした。

\subsection{実測 1：計算律速の処理は 4 コアでスケールする}

表 \ref{tab:rpi3-compute} に計算律速のベンチマークを示す。値は全て
安定状態（再起動後、クロックが落ち着いた状態、アクター未起動）での測定であり、
第\ref{sec:ratio}項のとおり比を主指標とする。

\begin{table}[htbp]
\centering
\caption{Pi 3：計算律速の速度向上（\route{/bench}、\texttt{cores\_online = 4}）}
\label{tab:rpi3-compute}
\small
\begin{tabular}{@{}llrrcc@{}}
\toprule
ベンチマーク & 性質 & 1 コア (\textmu s) & 4 コア (\textmu s) & 速度向上 & \texttt{agree} \\
\midrule
\bench{dining n=5}        & 純レジスタ演算   & 119{,}626  & 29{,}920  & \spd{3.99} & yes \\
\bench{nqueens n=12}      & 純計算（再帰）   & 260{,}053  & 70{,}149  & \spd{3.70} & yes \\
\bench{nqueens n=11}      & 純計算（再帰）   & 50{,}803   & 15{,}022  & \spd{3.38} & yes \\
\bench{nqueens n=13}      & 純計算（再帰）   & 1{,}419{,}536 & 445{,}639 & \spd{3.18} & yes \\
\bench{primes n=200000}   & 計算（下記参照） & 155{,}686  & 75{,}958  & \spd{2.04} & yes \\
\bottomrule
\end{tabular}
\end{table}

\noindent
\bench{dining} が理想値 \spd{4.00} にほぼ届く \spd{3.99} を示したことは、
本設計の妥当性を端的に示している。この処理はメモリに一切触れない純粋な
レジスタ演算であり、コア間で奪い合う資源が存在しないためである。

\bench{primes} だけが \spd{2.04} と低い。これは偶然ではなく、\textbf{仕事の
配分がこの問題の構造と共鳴した}結果である。ストライド 4 で配ると、コア 0 と
コア 2 は「4 の倍数」と「$4n+2$」すなわち\textbf{全て偶数}を担当し、
$d = 2$ で即座に合成数と判明して一瞬で終わる。一方コア 1 と 3 が奇数を
全て引き受ける。つまり実質 2 コアしか働いておらず、それがそのまま
\spd{2.04} という数字になっている。N-Queens では同じストライド配分が
正しく機能している（コストが列位置に対してなめらかで、偶奇と相関しないため）。
ブロック巡回配分に変えれば解消するが、これはベンチマークの均衡の問題であり、
本レポートの結論には影響しない。

\subsection{実測 2：メモリ律速の処理はスケールしない}
\label{sec:rpi3-fill}

表 \ref{tab:rpi3-fill} に、大量のメモリ書込（\bench{fill}）の測定を示す。
これは描画（フレームバッファへの書込）と同じ性質の処理であり、
ウィンドウシステムを並列化すべきかを判定するために用意した。

\begin{table}[htbp]
\centering
\caption{Pi 3：メモリ書込律速の速度向上（\bench{fill}、8 パス）}
\label{tab:rpi3-fill}
\small
\begin{tabular}{@{}rrrc@{}}
\toprule
語数 & 1 コア (\textmu s) & 4 コア (\textmu s) & 速度向上 \\
\midrule
64      & 5      & 13     & \spd{0.38} \\
256     & 11     & 18     & \spd{0.61} \\
1{,}024 & 42     & 37     & \spd{1.13} \\
4{,}096 & 174    & 136    & \spd{1.27} \\
16{,}384 & 689   & 511    & \spd{1.34} \\
65{,}536 & 2{,}665 & 2{,}137 & \spd{1.24} \\
262{,}144 (=1\,MB) & 11{,}007 & 8{,}125 & \spd{1.35} \\
\bottomrule
\end{tabular}
\end{table}

\noindent
2 つの重要な事実が読み取れる。

第一に、\textbf{どれだけ仕事を大きくしても \spd{1.35} 程度で頭打ち}になる。
1\,MB（1280$\times$800 の画面の約 1/4 に相当）を書いても \spd{1.35} であり、
計算律速の \spd{3.99} とは比較にならない。なお高クロック状態で測ると
同じ 1\,MB が \spd{0.95}、すなわち\textbf{4 コアの方がわずかに遅い}。
いずれの状態でも、並列化に意味がないという結論は変わらない。

第二に、\textbf{小さい仕事は並列化すると遅くなる}。64 語では \spd{0.38}、
すなわち 2.6 倍遅い。これはプールの往復コストに負けているためである。

\subsection{実測 3：プールの往復コストと損益分岐点}

何もしないジョブを投入して、プール自体の往復コストを測定した
（\bench{null}）。結果は\textbf{7--11\,\textmu s}であった。

これが並列化の\textbf{損益分岐点}である。すなわち、これを下回る仕事を
ワーカーに投げると必ず損をする。効率よく利得を得るには、その 10 倍、
すなわち\textbf{1 メッセージあたり 70\,\textmu s 以上の純計算}が必要となる。
この数値が、以降の適用判断の基準となった。

\subsection{適用判断 1：ウィンドウシステムは並列化しない}

「ウィンドウシステムなどのアプリケーションが 4 コアで速くなるか」という
問いに対し、上記の測定は\textbf{否}と答える。理由は 2 つあり、
どちらか一方だけでも致命的である。

\begin{enumerate}[leftmargin=1.6em]
  \item \textbf{描画はメモリ律速である。}このポートは D-cache が OFF なので、
        描画とはフレームバッファへの書込そのものであり、
        第\ref{sec:rpi3-fill}項の \bench{fill} と同じ性質を持つ。
        測定値は \spd{0.95}--\spd{1.35} であり、コアを足しても速くならない。
  \item \textbf{粒度が足りない。}ウィンドウシステムの描画 1 回は小さい。
        たとえばカーソルは 12$\times$12 = 144 画素であり、
        往復コスト 7--11\,\textmu s を下回る。実測でも 64 語の仕事は
        \spd{0.38} と\textbf{遅くなっている}。並列化すれば体感はむしろ悪化する。
\end{enumerate}

\noindent
したがってウィンドウシステムには手を加えなかった。
\textbf{速くならない上に遅くなる変更を入れないことも、工学上の判断である。}
このボードにおける法則は「\textbf{計算はスケールする、描画はスケールしない}」
と要約できる。

\subsection{適用判断 2：AIPL アクターは 1 箇所だけが該当する}

Pi 3 の AIPL ランタイムは、アクター 1 体につき Xinu スレッド 1 本と
メールボックス（リングバッファ＋計数セマフォ）を持ち、スレッドが
メッセージを取り出して \texttt{dispatch()} でメソッド本体を実行する
モデルである。

損益分岐点 70\,\textmu s に照らすと、\textbf{ほとんどの AIPL メソッドは
対象外}である。\texttt{bump} や ping-pong は算術数命令と \texttt{send} で
構成され、基準を 2--3 桁下回る。加えて \texttt{send}・\texttt{print}・
\texttt{wait} といったカーネル呼び出しは、そもそもワーカーコアで実行できない。

唯一の例外がロードバランサの \texttt{Worker.compute\_nq} である。

\begin{lstlisting}[caption={対象となる構造（\texttt{apps/abcl\_program.c}）}]
  int  ncores = smp_cores_online();
  unsigned long t0 = clkcount();
  long result = (long)abcl_nq_count_partial_smp((int)n, (int)first_col, ncores);
  long elapsed = (long)((clkcount() - t0) / 1000UL);   /* us -> ms */
  ...
  enqueue(self_id, disp, "done", 3, ...);   /* kernel call: stays on core 0 */
\end{lstlisting}

\noindent
重い計算が \texttt{abcl\_nq\_count\_partial()} の一点に隔離され、その前後が
カーネル呼び出しという理想的な構造である。n=13 で 1 メッセージあたり
約 100\,ms の純計算であり、基準を 3 桁上回る。ここだけを 4 コア化し、
カーネル呼び出しはコア 0 に残した。

\subsubsection*{分割の方法}

first\_col 1 つ分は\textbf{1 回の再帰呼び出し}であるため、分割は 1 段下、
すなわち\textbf{2 行目の合法な列}で行う。各コアにその剰余類を配る。
連続分割ではなくインターリーブとしたのは、コストが列位置に対して
中央が重く左右対称であり、連続分割では中央のコアに仕事が偏るためである。

\textbf{各コアは専用の盤面配列を持つ必要がある}。\texttt{nq\_recurse} は
探索の下降中に \texttt{cols[]} を破壊的に更新するため、配列を共有すると
コア同士が互いの探索を壊すからである。この配列を静的変数ではなく
呼び出し側スタックに置いたのは、プールが埋まって inline 実行に落ちた
スレッドと、プールを保持するスレッドとが、\textbf{どちらもコア 0 上で}
同じ配列を奪い合うためである。

ホスト上での総当たり検証により、この分割の正しさとバランスを確認した。

\begin{lstlisting}[language=,caption={カーネルのソースから当該関数を抜き出したホスト検証}]
REAL-SOURCE partial n=1..12 x every first_col x nc=1..4: ok
REAL-SOURCE summed partials vs known N-Queens counts: ok

balance n=13 first_col=6 nc=4 (per-core counts):  2233  2233  1802  1802
                                                        max/avg = 1.11
\end{lstlisting}

\subsubsection*{実機での結果}

アクター経路（\route{/api/loadbal/nqueens?n=12}）は正しく完走した。

\begin{lstlisting}[language=,caption={アクター経由の N-Queens（n=12 の正解は 14200）}]
nqueens n=12 submitted=12 task_id_range=1..12
nqueens-result done=1 received=12 expected=12 total=14200
\end{lstlisting}

\noindent
答えは正しく、また各 Worker の \texttt{busy\_ms} に\textbf{値が入るようになった}
（従来は壊れた \texttt{clkticks} を使っていたため常に 0 であった）。

\subsection{発見した不具合}
\label{sec:rpi3-bugs}

本作業では、実機に投入する前のホスト検証と、測定データの精査により、
以下の不具合を発見・修正した。いずれも\textbf{測定していなければ
気付けなかった}ものである。

\begin{enumerate}[leftmargin=1.6em]
  \item \textbf{素数分割の剰余類の誤り（自作、投入前に発見）。}
        開始値の算出を誤り、剰余類がずれて一部の値を取りこぼし、
        かつストライド 1 のとき 1 を素数に数えていた。ホスト検証で捕捉。

  \item \textbf{ジョブ引数のスレッド間競合（自作、投入前に発見）。}
        第\ref{sec:ud}項に述べた競合。アクターは Worker が 4 体、
        すなわちスレッド 4 本であり、実際に起こり得た。API に \texttt{ud} を
        追加して構造的に解消した。

  \item \textbf{\texttt{smp\_parallel\_sum} の再入不可（自作、投入前に発見）。}
        単一のジョブメールボックスを複数スレッドが同時に使うと破壊し合う。
        使用中なら呼び出し側で全チャンクを実行する方式で解消した。

  \item \textbf{タイマ割り込み由来の時計が常に 0（既存）。}
        第\ref{sec:clock}項。本作業以前から存在した。
        \texttt{clkcount()} 直読で回避した。
\end{enumerate}

\subsection{副次的な発見：\texttt{volatile} グローバルの代償}
\label{sec:rpi3-volatile}

第\ref{sec:ud}項の競合対策として、ジョブ引数を \texttt{volatile} な
ファイルスコープ変数から呼び出し側スタックの構造体へ移したところ、
\textbf{速度向上比が \spd{3.04} から \spd{3.70} へ改善した}。

理由は D-cache が OFF であることにある。\texttt{volatile} 変数は参照のたびに
必ず実メモリを読むため、D-cache のないこの環境では\textbf{内側ループの
1 反復ごとに RAM 往復}が発生していた。通常の構造体メンバに変えたことで
コンパイラがレジスタに載せられるようになった。

すなわち\textbf{競合バグの修正が、同時に性能改善でもあった}。
D-cache を持たない環境では、\texttt{volatile} の乱用は通常環境より遥かに高くつく。


% =======================================================================
% ボード追加時はここの コメントを外す。
% \section{Raspberry Pi 4 (BCM2711)}
% \label{sec:rpi4}
% % ======================================================================
% boards/rpi4.tex — Raspberry Pi 4 (BCM2711, Cortex-A72) のマルチコア化
% 実測: 2026-07-19, HTTP /bench (192.168.3.100), production kernel md5 4cf1a50a
% ======================================================================

Raspberry Pi 4（BCM2711、Cortex-A72 $\times$ 4、AArch64）へ、第 \ref{sec:design} 節と
同一のワーカープール型 SMP を移植した。共通設計・測定方法はボードに依存しないので、
ここでは\textbf{Pi 4 固有の差分と実測値}のみを述べる。

\subsection{対象と前提}

BCM2711 は 4 つの Cortex-A72 を持つ。Pi 3（A53、32bit）と異なり\textbf{AArch64
（64bit）}で動作する。ファームウェアはコアを EL2 に落とすため \texttt{boot.S} で
EL2$\to$EL1 に遷移させる。D-cache は第 \ref{sec:design} 節の前提どおり\textbf{OFF}
（\texttt{SCTLR.C=0}、実測 \texttt{SCTLR}=\texttt{0x30d01801}）であり、ロックフリー
調停が RAM 直行で成立する。この D-cache を敢えて ON にする独立実験は第
\ref{sec:dcache-a72} 節に示した。

\subsection{コア起動：spin table}
\label{sec:rpi4-bringup}

Pi 4 のファームウェア（armstub8）はコア 1--3 を物理アドレス
\texttt{0xe0}/\texttt{0xe8}/\texttt{0xf0} でスピンさせて保持する。各スロットへ
エントリ関数アドレスを書き \texttt{SEV} で起こす（Pi 3 の mailbox、Pi 5 の PSCI と
対応する差分）。起床した各コアは \texttt{mmu\_enable\_secondary} で MMU＋I-cache を
コア 0 と同一構成にし、\texttt{cores\_online = 4} を実機で確認した。

\subsection{実測 1：計算律速はスケール、ただし配分方式に強く依存}

\begin{table}[htbp]
\centering
\caption{Pi 4：計算律速のコア数スイープ（\route{/bench}、\texttt{cores=1..4}、全て \texttt{agree=yes}）}
\label{tab:rpi4-compute}
\small
\begin{tabular}{@{}lccccl@{}}
\toprule
ベンチ & 1 & 2 & 3 & 4 & 備考 \\
\midrule
\bench{dining n=5}      & \spd{0.99} & \spd{1.99} & \spd{3.00} & \spd{3.99} & 完全に線形（理想） \\
\bench{primes n=300000} & \spd{0.99} & \spd{1.61} & \spd{2.30} & \spd{3.00} & 連続分割で概ね均衡 \\
\bench{nqueens n=12}    & \spd{0.99} & \spd{1.60} & \spd{1.79} & \spd{1.85} & 静的分割では頭打ち \\
\bottomrule
\end{tabular}
\end{table}

\noindent
\bench{dining}（独立テーブルの均一コスト）は Pi 3・Pi 5 と同じく理想の
\spd{3.99} を示す。\bench{nqueens} は 4 コアで \spd{1.85}——だがこれは
分割方式（連続 vs インターリーブ）の問題ではなく\textbf{静的分割そのものの限界}で
あることを、第 \ref{sec:cross-fair} 項の実験（連続・第 1 行・第 2 行の三方式が
すべて \spd{2.0} 前後）で確認した。動的負荷分散を要する Level 3 の対象である。

\subsection{実測 2：メモリ律速はスケールしない}

\begin{table}[htbp]
\centering
\caption{Pi 4：メモリ書込律速（\bench{fill}、8 パス、\texttt{cores=4}）}
\label{tab:rpi4-fill}
\small
\begin{tabular}{@{}rrrl@{}}
\toprule
語数 & 1 コア \textmu s & 4 コア \textmu s & 速度向上 \\
\midrule
64      & 3     & 7     & \spd{0.42} \\
256     & 8     & 12    & \spd{0.66} \\
1{,}024 & 28    & 35    & \spd{0.80} \\
4{,}096 & 110   & 100   & \spd{1.10} \\
16{,}384 & 438  & 307   & \spd{1.42} \\
65{,}536 & 1{,}749 & 1{,}269 & \spd{1.37} \\
262{,}144 (=1\,MB) & 6{,}997 & 5{,}050 & \spd{1.38} \\
\bottomrule
\end{tabular}
\end{table}

\noindent
Pi 3 と同じ二つの現象が現れる：\textbf{大きな仕事でも \spd{1.4} で頭打ち}（1\,MB で
\spd{1.38}）、\textbf{小さな仕事は並列化すると遅くなる}（64 語で \spd{0.42}）。
D-cache OFF ゆえメモリ書込が RAM 帯域律速になるためで、計算律速の \spd{3.99} とは
比較にならない。

\subsection{実測 3：プールの往復コスト}

空ジョブ（\bench{null}）の 4 コア投入は実測\textbf{約 5\,\textmu s}である（1 コアは
0\,\textmu s）。この往復コストを下回る仕事は並列化すると損をする（表
\ref{tab:rpi4-fill} の 64 語 \spd{0.42} が実例）。Pi 3 の 7--11\,\textmu s より
小さいのは A72 のクロック・キャッシュがコア間シグナリングを速めたためである。

\subsection{小括}

Pi 4（A72）は、\textbf{計算律速は理想的にスケール（\spd{3.99}）・メモリ律速は
スケールしない（\spd{1.38}）}という Pi 3 の結論をそのまま再現した。相違は絶対速度
（A72 は A53 より速い）と、プール往復の小ささ、そして AArch64／spin table という
起動機構のみである。「並列化すべき対象を分ける境界」は Pi 3 と同一であった。

%
% \section{Raspberry Pi 5 (BCM2712)}
% \label{sec:rpi5}
% % ======================================================================
% boards/rpi5.tex — Raspberry Pi 5 (BCM2712, Cortex-A76) のマルチコア化
% 実測: 2026-07-19, HTTP /bench (192.168.3.101)
% ======================================================================

Raspberry Pi 5（BCM2712、Cortex-A76 $\times$ 4、AArch64）へも同一のワーカープール型
SMP を移植した。ここでも\textbf{Pi 5 固有の差分と実測値}のみを述べる。

\subsection{対象と前提}

BCM2712 は 4 つの Cortex-A76 を持つ。A76 は DynamIQ Shared Unit（DSU）を備え、
D-cache ON 時のコア間コヒーレンシ機構が A53/A72 の \texttt{SMPEN} とは異なる（第
\ref{sec:dcache-a76} 節）。本節の測定は本レポート共通の前提どおり\textbf{D-cache
OFF}で行い、ロックフリー調停を RAM 直行で成立させている。A76 は 3 世代で\textbf{単一
コアの絶対速度が最も速い}（例：\bench{fill} 1\,MB の 1 コアが約 2{,}170\,\textmu s で、
A72 の約 7{,}000\,\textmu s の 3 倍以上速い）。

\subsection{コア起動：PSCI}
\label{sec:rpi5-bringup}

Pi 5 はコア起動に\textbf{PSCI}（\texttt{CPU\_ON} のハイパーコール）を用いる（Pi 3 の
armstub mailbox、Pi 4 の spin table に対応する差分）。コア識別は \texttt{MPIDR} の
Aff1 フィールドで行う。起床後の MMU 構成統一は他ボードと同一で、
\texttt{cores\_online = 4} を実機で確認した。

\subsection{実測 1：計算律速}

\begin{table}[htbp]
\centering
\caption{Pi 5：計算律速のコア数スイープ（\route{/bench}、\texttt{cores=1..4}、全て \texttt{agree=yes}）}
\label{tab:rpi5-compute}
\small
\begin{tabular}{@{}lccccl@{}}
\toprule
ベンチ & 1 & 2 & 3 & 4 & 備考 \\
\midrule
\bench{dining n=5}      & \spd{1.00} & \spd{1.99} & \spd{2.99} & \spd{3.99} & 完全に線形（理想） \\
\bench{primes n=300000} & \spd{1.00} & \spd{1.61} & \spd{2.31} & \spd{3.01} & 連続分割で概ね均衡 \\
\bench{nqueens n=12}    & \spd{1.00} & \spd{1.45} & \spd{1.47} & \spd{1.5}\,--\spd{1.5} & 静的分割では頭打ち \\
\bottomrule
\end{tabular}
\end{table}

\noindent
\bench{dining} は A53・A72 と同じく理想の \spd{3.99}。\bench{nqueens} は 4 コアで
\spd{1.5} 前後（\spd{1.41}--\spd{1.54}）と、\textbf{A72（\spd{1.85}）よりさらに低い}。
これは第 \ref{sec:cross-fair} 項の知見と整合する：A76 は単一コアが速いぶん、静的分割の
負荷不均衡とプール往復の相対比重が大きくなり、\bench{nqueens} の静的分割はより早く
頭打ちになる。改善には動的負荷分散（Level 3）が要る。

\subsection{実測 2：メモリ律速は最も平坦}

\begin{table}[htbp]
\centering
\caption{Pi 5：メモリ書込律速（\bench{fill}、8 パス、\texttt{cores=4}、5 回中の最良値）}
\label{tab:rpi5-fill}
\small
\begin{tabular}{@{}rrrl@{}}
\toprule
語数 & 1 コア \textmu s & 4 コア \textmu s & 速度向上 \\
\midrule
4{,}096  & 34    & 31    & \spd{1.09} \\
16{,}384 & 187   & 183   & \spd{1.02} \\
65{,}536 & 536   & 518   & \spd{1.09} \\
262{,}144 (=1\,MB) & 2{,}169 & 2{,}043 & \spd{1.06} \\
\bottomrule
\end{tabular}
\end{table}

\noindent
Pi 5 の \bench{fill} は全域で \spd{1.0}--\spd{1.1} と、3 世代で\textbf{最も平坦}で
ある。A76 は 1 コアあたりのメモリ帯域が相対的に高く、\textbf{1 コアで既に RAM 帯域を
ほぼ飽和}させるため、コアを足しても増えない。なお \bench{fill} の小サイズ（$\le$1{,}024
語、サブ\textmu s）は約 8 秒周期の自動アクタ GC に撹乱されて外れ値が出るため、5 回
測定の最良値で代表させた——単発測定のばらつき（第 \ref{sec:cross-fair} 項でも触れた）
の一例である。

\subsection{小括}

Pi 5（A76）でも本レポートの中心結論はそのまま成立する：\textbf{計算律速は理想的に
スケール（\spd{3.99}）・メモリ律速はスケールしない（\spd{1.06}）}。世代固有の差は、
起動機構（PSCI）、単一コア絶対速度が最速であること、\bench{fill} が最も平坦である
こと、そして D-cache ON 時のコヒーレンシ機構が DSU である点（第 \ref{sec:dcache-a76}
節）に限られる。「並列化すべき対象を分ける境界」は 3 世代を通じて同一であった。


% =======================================================================
\section{実験：D-cache を有効化する（Level 2）}
\label{sec:dcache}
% ======================================================================
% experiments/dcache.tex — Level 2: enabling the D-cache
%
% ★ 軸はボードではなく「キャッシュ構成」。ゆえに boards/ ではなく
%    experiments/ に置く。同じ枠組みで別の軸を切りたくなったら
%    experiments/ に足す（例: 動的ワークスティーリング = Level 3）。
%
%    設計・制約・DMA棚卸し・ハーネスは測定に依存せず確定。
%    D-cache ON の「スケールするか」の実測値だけがシリアル取得待ち
%    （§末尾の \TODO を、数値が出たら差し替える）。
% ======================================================================

\subsection{動機：支配変数はキャッシュ構成である}

第 \ref{sec:rpi3} 節の測定で最も効いていた制約は、SMP の方式ではなく
\textbf{D-cache が OFF であること}であった。メモリ律速がスケールしなかったのも
（\bench{fill} \spd{0.93}）、\texttt{volatile} を外したら \spd{3.04} が
\spd{3.66} に上がったのも（第 \ref{sec:rpi3-volatile} 項）、動的ワークスティーリングを
断念して静的分割にしたのも、すべて D-cache OFF が原因である。

そこで問いは：\textbf{D-cache を ON にすると、メモリ律速の処理はスケールし始めるか}。
もしそうなら「ウィンドウシステムは並列化しない」という第 \ref{sec:rpi3}\,節の
結論が覆る。この問いは SMP の方式とは独立であり、キャッシュ構成という別の軸に属する。

\subsection{D-cache ON は「フラグを 1 にする」ではない}

素朴に \texttt{SCTLR.C} を 1 にするのは\textbf{危険で、過去に一度 Pi 3 を
brick させている}（\texttt{mmu.c} 冒頭に記録あり）。安全に ON にするには、
この移植に 1 行も存在しなかったキャッシュ管理を新規に書く必要があった。
本実験で顕在化した罠は 4 つある。

\begin{enumerate}[leftmargin=1.6em]
  \item \textbf{リセット直後のキャッシュはゼロでなくゴミ。}
        invalidate せずに \texttt{C=1} にすると、コアはゴミを有効データと
        見なす。これが brick の原因である。set/way invalidate ループを
        \texttt{C=1} の\textbf{前に}実行しなければならない（後では、いま
        キャッシュした本物のデータを捨ててしまう）。

  \item \textbf{Cortex-A53 は SMPEN なしでは非コヒーレント。}
        \texttt{mmu.c} は RAM を shareable にマップしているが、
        A53 では shareable かつ cacheable なメモリでも \texttt{CPUECTLR.SMPEN}
        を立てるまで\textbf{コア間でコヒーレントにならない}。\texttt{C=1} かつ
        SMPEN=0 だと、第 \ref{sec:design} 節のロックフリー・ジョブメールボックス
        （\texttt{volatile} + \texttt{dsb}）が静かに壊れる。\texttt{volatile} は
        コンパイラを縛るだけで、キャッシュは縛らない。

  \item \textbf{二次コアの起動窓。}
        二次コアは \texttt{\_smp\_start} の時点で\textbf{まだ MMU も
        キャッシュも OFF}であり、その状態でスタックポインタを RAM から読む。
        この読みはコア 0 のキャッシュをスヌープしない。ゆえにコア 0 は
        スタックアドレスを\textbf{起こす前に clean}せねばならず、さもなくば
        二次コアは null スタックで起きて最初の push で死ぬ。

  \item \textbf{\texttt{mmu\_disable} の前提が崩れる。}
        「キャッシュ OFF だから flush 不要」というコメントが \texttt{C=1} で
        偽になる。\texttt{kexec} 前にキャッシュを落とすと、\texttt{/upload} が
        キャッシュ経由で書いた次カーネルのイメージが dirty ライン内に取り残され、
        \textbf{stale メモリへジャンプ}する。clean を追加した。
\end{enumerate}

\noindent
これらを \texttt{cache.c}（set/way・範囲・SMPEN の各操作を新規実装）と、
\texttt{mmu.c} の \texttt{DCACHE\_ON} 経路（順序が命：invalidate と SMPEN を
\texttt{C=1} の前に）で処理した。

\subsection{最大の障害：USB DMA はページ属性で直せない}
\label{sec:dcache-usb}

D-cache ON でも DMA するハードとの共有メモリは非コヒーレントになる。
そこでカーネル全体の\textbf{DMA 地点を棚卸し}した。結果は明快で、
FB・電源・WiFi のメールボックスや SD/SDIO は\textbf{全て PIO}（DATA レジスタ
経由）であり、コヒーレンシ問題を持たない。唯一の DMA エンジンは
\textbf{DWC USB}であり、これが有線 LAN を担っている。

そして USB こそが\textbf{構造的に直せない}地点であった。DWC ドライバは
呼び出し側のポインタをそのまま DMA アドレスに渡し、呼び出し側は
\textbf{ヒープ}バッファ（\texttt{memget} は 8 バイト境界＝隣の割り当てと
64 バイトのキャッシュラインを共有）と、あろうことか\textbf{スタック上の
局所変数}（\texttt{smsc9512.c} のレジスタ操作が 4 バイトを DMA する）を渡す。
64 バイトのラインを invalidate すれば、\textbf{同じラインにある戻り番地や
退避レジスタごと破棄}する。これはページ属性では解けず、バウンスバッファと
ドライバ全域の 64 バイト境界化を要する。

この発見自体が結論である：\textbf{D-cache を入れられないのはキャッシュ管理を
書くのが面倒だからではなく、ドライバの DMA バッファの出自が構造的に許さない
から}である。

\subsection{実験の切り分け：USB を解かずに問いに答える}

問いは「D-cache でメモリ律速はスケールするか」であり、これに答えるのに
USB を解く必要はない。そこで実験用ビルド変種を用意した。

\begin{center}
\texttt{-DDCACHE\_EXPERIMENT -DDCACHE\_ON -DOS\_MINIMAL}
\end{center}

\noindent
\texttt{OS\_MINIMAL} でネットと WM を落とし、\texttt{DCACHE\_EXPERIMENT} で
さらに \texttt{usbinit}（唯一の DMA エンジン）を落とす。これで DMA する相手が
一つも存在しない状態で \texttt{C=1} にでき、USB の構造問題を回避したまま
本題を測れる。ただし\textbf{HDMI は残した}——ユーザ要求であり、また
フレームバッファの扱いが D-cache ON 実装の良い試金石でもあるためである。

HDMI を残すには 2 つの手当てが要った。FB のメールボックス要求構造体は
元は\textbf{スタック上}にあり（GPU が stale RAM を読み \texttt{address=0}
＝SYSERR ＝画面出ず）、64 バイト境界・64 バイト詰めの専用 static に移して
handshake の前後で clean/invalidate した。FB のピクセル領域自体は、
GPU がアドレスを返した\textbf{後に} \texttt{mmu\_set\_range\_noncached()} で
Normal Non-cacheable へ貼り替えた（アドレスは \texttt{mmu\_init} の後に
判明するため、静的な表には焼けない）。

\subsection{測定ハーネス}

実験変種はネットが無いので \route{/bench} を提供できない。そこでワークロード本体を
\texttt{apps/smpbench.c} に切り出し、\textbf{HTTP 経路（既定カーネル）と起動時
コンソール経路（実験カーネル）が同一のコードを呼ぶ}ようにした。キャッシュ
ON/OFF の比較は、両者が同一実装を走らせて初めて意味を持つ——別実装なら
「キャッシュ設定の差」ではなく「実装の差」を測ってしまう。この切り出しの
過程で、第 \ref{sec:rpi3-measbug} 項の測定バグを発見・修正している。

実験カーネルは起動時に \texttt{smpbench\_run\_all()} を実行し、\bench{fill} の
サイズ掃引を中心とした結果をシリアルへ出力する。

\subsection{実測結果}
\label{sec:dcache-result}

D-cache ON カーネルを実機に投入し、起動時ハーネスの出力をシリアルで取得した。
まず 3 つの前提が\textbf{すべて確認された}。

\begin{itemize}[leftmargin=1.6em]
  \item \textbf{起動した（brick 回避）}。Xinu のバナーからシェルまで正常に立ち
        上がった。set/way invalidate を \texttt{C=1} の前に置く順序が効いている。
  \item \textbf{\texttt{cores\_online = 4}、全ベンチ \texttt{agree = ok}}。
        SMPEN を立てたことで、\texttt{C=1} 下でもロックフリー・ジョブメールボックスの
        コヒーレンシが保たれている。逐次と並列の答えは全ベンチで一致した。
  \item \textbf{\texttt{sctlr = 0x00c5183d}}。D-cache OFF の \texttt{0x00c51839} に
        対しビット 2（C）が立ち、実機で D-cache が有効化されたことを裏付ける。
\end{itemize}

\subsubsection*{本命：メモリ律速はキャッシュでスケールする——ただし容量まで}

\begin{table}[htbp]
\centering
\caption{\bench{fill} の速度向上：D-cache OFF vs ON（8 パス、作業セット掃引）}
\label{tab:dcache-fill}
\small
\begin{tabular}{@{}rccl@{}}
\toprule
語数（作業セット） & OFF の速度向上 & ON の速度向上 & \\
\midrule
64          & \spd{0.38} & \spd{2.00} & \\
256         & \spd{0.61} & \spd{2.50} & \\
1{,}024 (4\,KB)   & \spd{1.13} & \spd{3.60} & \\
4{,}096 (16\,KB)  & \spd{1.27} & \spd{3.88} & \\
16{,}384 (64\,KB) & \spd{1.34} & \spd{4.01} & ← 完全スケール \\
65{,}536 (256\,KB)& \spd{1.24} & \spd{2.51} & ← 低下開始 \\
262{,}144 (1\,MB) & \spd{1.35} & \spd{1.40} & ← 帯域律速へ回帰 \\
\bottomrule
\end{tabular}
\end{table}

\noindent
表 \ref{tab:dcache-fill} が本実験の答えである。第 \ref{sec:rpi3} 節の結論
「メモリ律速はスケールしない」は、\textbf{D-cache OFF に固有の現象}であった。
D-cache ON では、作業セットがキャッシュに収まる限り\textbf{メモリ書込は
理想的にスケールする}（64\,KB で \spd{4.01}）。

しかしスケールするのは容量までである。作業セットが 256\,KB を超えると速度向上は
崩れ、1\,MB では \spd{1.40} まで戻る。これは\textbf{キャッシュ容量効果}である。
小さい \bench{fill} は同じ領域を 8 回書くため 2 回目以降はキャッシュに当たり、
実質\textbf{計算律速}になってスケールする。作業セットがキャッシュ（Pi 3 の
Cortex-A53 は L1 データ 32\,KB／コア＋共有 L2 512\,KB）を超えると、
書き戻しと充填で\textbf{再び帯域律速}に戻る。速度向上が 64\,KB で頂点、
256\,KB で崩れ始めるのは、この容量境界と整合する。

\subsubsection*{ウィンドウシステムの結論は覆るか——覆らない}

一見、これは「WM は並列化しない」（第 \ref{sec:rpi3} 節）を覆すように見える。
だが\textbf{覆らない}。理由が 2 つある。

第一に、\textbf{実フレームバッファはキャッシュ不可にマップしてある}
（第 \ref{sec:dcache-usb} 項、\texttt{mmu\_set\_range\_noncached}）。GPU が
スヌープしないため、これは D-cache ON でも変えられない。表 \ref{tab:dcache-fill}
の \bench{fill} は\textbf{通常のキャッシュ可能バッファ}での測定であり、
実描画はこの恩恵を受けない。

第二に、恩恵を受けるのは作業セットがキャッシュに収まる場合だけだが、
全画面ブリットは 1280$\times$800$\times$4 = 3.9\,MB でキャッシュを大きく超え、
仮にキャッシュ可能でも \spd{1.40} 程度に留まる。

したがって結論は次のように\textbf{精緻化}される：メモリ律速の処理は、
D-cache ON かつ\textbf{作業セットがキャッシュ内}なら 4 コアでスケールする。
だが実描画は（キャッシュ不可の FB へ、キャッシュを超える量を書くため）
どちらの条件も満たさず、依然スケールしない。中間バッファ上での
合成のように、キャッシュ内で完結する処理を挟めば話は別である。

\subsubsection*{副産物：D-cache は単一コア性能も大きく上げる}

計算律速側では、\bench{dining}(\spd{3.99})・\bench{primes}(\spd{2.05}) の
速度向上比は D-cache OFF とほぼ不変であった（元々スケールしていたため）。
一方で\textbf{絶対性能}は激変した。\bench{nqueens n=12} の 1 コア時間は
D-cache OFF の約 167{,}000\,\textmu s から\textbf{約 19{,}000\,\textmu s へ、
実に 8.8 倍高速化}している。クロック差（起動直後 2.33 倍）だけでは説明できず、
再帰探索が触る盤面配列 \texttt{cols[]} とスタックがキャッシュに載った効果である。
並列スケーリングとは独立に、D-cache は\textbf{計算律速の絶対性能そのもの}を
底上げする。

\subsubsection*{ただし本番投入はできない}

以上は\textbf{実験変種}（\texttt{DCACHE\_EXPERIMENT}、USB を落とした構成）での
結果である。本番カーネルに D-cache を入れるには、第 \ref{sec:dcache-usb} 項の
USB DMA 問題——ヒープとスタックのバッファへのバウンス、およびドライバ全域の
64 バイト境界化——を解く必要がある。それは本実験の範囲外であり、次段階の課題である。

\subsection{付録：実機シリアル出力（D-cache ON）}

参考のため、実機（コールドブート）から取得した起動時ハーネスの生出力を示す。

\begin{lstlisting}[language=,basicstyle=\ttfamily\scriptsize]
===== smpbench =====
cores_online = 4
D-cache      = ON   (SCTLR.C=1, SMPEN set)
sctlr        = 0x00c5183d

-- compute-bound (expect x3.2-x4.0) --
  bench           1-core us  N-core us  speedup
  dining n=5          51290      12828   x3.99  ok
  nqueens n=11         3716       1027   x3.61  ok
  nqueens n=12        19024       4865   x3.91  ok
  primes 200k         66761      32550   x2.05  ok

-- memory-store-bound: THE QUESTION (x1.0 = no benefit) --
  words           1-core us  N-core us  speedup
  fill 64                 2          1   x2.00  ok
  fill 256                5          2   x2.50  ok
  fill 1024              18          5   x3.60  ok
  fill 4096              70         18   x3.88  ok
  fill 16384            285         71   x4.01  ok
  fill 65536           1177        468   x2.51  ok
  fill 262144          4682       3330   x1.40  ok

-- pool round trip (the parallelise/don't break-even) --
  null                    0          1   x0.00  ok

===== end smpbench =====
\end{lstlisting}
\subsection{本番化の試みと限界：USB DMA は越えられなかった}
\label{sec:dcache-prod}

第 \ref{sec:dcache-usb} 項で「USB は構造的に難しい」と述べたが、実際に
本番化を試みた。結果は\textbf{部分成功}であり、その失敗の様相そのものが
価値ある結果なので記録する。

\subsubsection*{方式：全 USB 転送を非キャッシュ・バウンス経由に}

USB DMA の障害は、バッファがヒープ（\texttt{memget} は 8 バイト境界＝
64 バイトのキャッシュラインを隣と共有）やスタック上にあり、per-transfer の
clean/invalidate が\textbf{隣接データを巻き込む}ことだった。そこで
per-transfer 管理を避け、\textbf{全転送を専用の非キャッシュ・アリーナ経由に
バウンス}する方式を採った。DWC ドライバは元々、非境界バッファ用のバウンス
経路（\texttt{aligned\_bufs} + TX/RX コピー）を持つので、その切替を
「D-cache ON なら常にバウンス」に変え、アリーナを非キャッシュにマップした。
非キャッシュ領域は CPU も DMA も同一 RAM を直接見るので、原理上キャッシュ
管理命令が一切不要になる。

\subsubsection*{実機結果：起動と HDMI は成功、USB は全滅}

\begin{table}[htbp]
\centering
\caption{D-cache 本番化（フルカーネル）の実機結果}
\small
\begin{tabular}{@{}ll@{}}
\toprule
項目 & 結果 \\
\midrule
D-cache ON で起動 & $\checkmark$ brick なし（デスクトップ表示） \\
HDMI / フレームバッファ & $\checkmark$ 動作（FB の非キャッシュ再マップが有効） \\
set/way invalidate + SMPEN & $\checkmark$ 機能 \\
単一コア性能 & $\checkmark$ 8.8 倍（実験変種で実測） \\
\textbf{USB（マウス・キーボード・ethernet）} & $\times$ 全滅 \\
\bottomrule
\end{tabular}
\end{table}

\noindent
HDMI が正しく映ることは、非キャッシュ化の機構そのもの
（\texttt{mmu\_set\_range\_noncached}）が\textbf{正しく動作している証拠}である
（フレームバッファも同じ関数で非キャッシュ化しており、壊れていれば画面が
乱れる）。にもかかわらず USB は全滅した。Raspberry Pi 3 では
マウス・キーボード・ethernet がすべて LAN9514 ハブの背後にあるため、
USB が一つ壊れると全滅する。実機シリアルログには
\texttt{[webactor] autostart: ETH0 never came up} が記録された。

\subsubsection*{四つの修正を試みた——いずれも解決せず}

原因を、実機を最小限しか使わずオフラインのコードレビューで追った。

\begin{enumerate}[leftmargin=1.6em]
  \item \textbf{DSB バリア追加。} 非キャッシュ書き込みもライトバッファに
        滞留し得るため、DMA 起動前に DSB で RAM へ排出。\textbf{必要だが
        不十分}——単独では USB は復活しなかった。

  \item \textbf{\texttt{packet\_count} オーバーフロー修正。} \texttt{packet\_count}
        は 10 ビット（最大 1023）で、バウンスアリーナを 128\,KB に拡大した際に、
        小さい \texttt{max\_packet\_size} で桁溢れし得ることを発見。しかし
        \textbf{赤鰊だった}——スプリット転送では \texttt{transfer.size} が
        \texttt{max\_packet\_size} に事前制限され \texttt{packet\_count} は
        常に 1 で溢れない。撤回した。

  \item \textbf{スプリット転送ロジックの精査。} Pi 3 は全転送がスプリット
        （ハブ背後）なので、バウンスとスプリット再試行の相互作用を徹底
        トレースした。\textbf{ロジックは正しかった}——\texttt{dma\_address} の
        保持、再アーム、コピーバックのオフセットはすべて整合していた。

  \item \textbf{アリーナ非キャッシュ化の SMP タイミング修正。} 最有力仮説。
        非キャッシュ化を実行時（\texttt{usbinit} 内）に、しかも現在コアだけ
        TLB 無効化で行っていたが、SMP のワーカーコアはそれ\textbf{以前}に
        アリーナがキャッシュ可能なまま起動しており、コア間でマッピングが
        不整合だった。アリーナのアドレスはコンパイル時に確定するので、
        \texttt{mmu\_init} で全コア起動前・D-cache 有効化前に非キャッシュを
        ベイクするよう変更した。しかし\textbf{これも USB を復活させなかった}。
\end{enumerate}

\subsubsection*{結論：本番化は USB DMA に阻まれた}

四つの修正を尽くしても USB は動かなかった。バウンスロジックとスプリット
処理は精査済みで正しく、コヒーレンシ機構も（HDMI が証明する通り）機能して
いる。それでも USB が全滅する——真因は、これらより深い、実機レベルの精密
診断（USB 列挙の詳細なシリアルログ等）を要する相互作用にある。手元の
シリアルアダプタ（Prolific、macOS で不安定）ではそこまで追えなかった。

したがって本レポートの範囲では、\textbf{D-cache は計算と表示には本番投入
できるが、USB DMA が最後の、そして越えられない壁である}と結論する。第
\ref{sec:dcache-usb} 項の構造的予測——USB のバッファ出自が本質的な障害——が、
バウンスという迂回策すら退ける形で、より強く裏付けられた。8.8 倍の単一コア
高速化という果実は実在するが、それを本番で収穫するには、USB ドライバの
DMA 経路をより根本的に作り直す必要がある。これは明確に次段階の課題である。


% =======================================================================
\section{ボード横断比較}
\label{sec:cross}

\begin{table}[htbp]
\centering
\caption{ボード別 ワーカープール型 SMP の比較（rpi5 は追記予定）}
\label{tab:cross}
\small
\begin{tabularx}{\textwidth}{@{}l l l l X@{}}
\toprule
 & \textbf{Pi 3 (BCM2837)} & \textbf{Pi 4 (BCM2711)} & \textbf{Pi 5 (BCM2712)} & \textbf{備考} \\
\midrule
CPU        & Cortex-A53 $\times$4 & Cortex-A72 $\times$4 & Cortex-A76 $\times$4 & \\
実行状態   & AArch32 (32bit)      & AArch64 & AArch64 & Pi 3 のみ 32bit \\
コア解放   & armstub7 mailbox 3   & spin table & PSCI & 第\ref{sec:rpi3-bringup}項 \\
D-cache    & OFF                  & OFF & OFF & ロックフリーの前提 \\
計算 \bench{dining} & \spd{3.99}  & \spd{3.99} & --- & 計算律速（公平比較） \\
メモリ \bench{fill} 1\,MB & \spd{0.9}--\spd{1.35} & \spd{1.40} & --- & メモリ律速（公平比較） \\
プール往復 & 7--11\,\textmu s      & 4\,\textmu s & --- & 並列化の損益分岐 \\
\bottomrule
\end{tabularx}
\end{table}

\subsection{公平な比較の条件：同一の分割方式のみ}
\label{sec:cross-fair}

Pi 4（A72）でも同一の \route{/bench} ハーネス（同一ワークロード・
\texttt{agree} 検証・\textmu s タイミング・\texttt{cores} スイープ）で測定した。
ただし\textbf{全ベンチが比較可能なわけではない}。Pi 3 と Pi 4 は N-Queens と
素数の\textbf{分割方式が異なる}——Pi 3 は本作業で実装したインターリーブ
（ストライド）分割、Pi 4 は既存の連続分割である。第 \ref{sec:rpi3-measbug} 項で
述べたとおり、分割方式は偏った負荷の速度向上を大きく左右する。実際、
\bench{nqueens} は Pi 3 \spd{3.65} に対し Pi 4 \spd{2.05}——だがこれは
微アーキの差ではなく、連続分割が N-Queens の中央偏重コストで負荷を偏らせた
結果である（第 \ref{sec:rpi3-fill} 項の max/avg 1.48 対 1.18 と整合）。
したがって\textbf{両ボードで分割方式が同一のベンチだけ}が公平な比較となる：
均一コストの \bench{dining}（両方連続）と \bench{fill}（両方連続）である。

\subsection{知見は微アーキを超えて一般化する}

その 2 つで、本レポートの中心的結論が A53 から A72 へ\textbf{一般化する}ことが
確認できた。

\begin{itemize}[leftmargin=1.6em]
  \item \textbf{計算律速は両ボードで理想的にスケールする。}\bench{dining} の
        コア数スイープは A53・A72 とも
        \spd{1.00}/\spd{1.99}/\spd{2.99}/\spd{3.99} と\textbf{同一の完全な直線}を
        描く。ワーカープール型 SMP による計算の線形スケーリングは、
        マイクロアーキテクチャに依存しない。
  \item \textbf{メモリ律速は両ボードでスケールしない。}\bench{fill} は A53 で
        \spd{0.9}--\spd{1.35}、A72 で全サイズ \spd{1.40} 程度と、どちらも
        \textbf{平坦}で線形から遠い。D-cache OFF ゆえメモリ書込が RAM 帯域律速に
        なる現象は、両ボード共通である（A72 の方がわずかに高いのは、より広い
        メモリサブシステムがコア間の重なりをやや許すためと考えられる）。
\end{itemize}

\noindent
すなわち「計算はスケールし、描画（メモリ）はスケールしない」という Pi 3 の
結論は、Pi 4 でもそのまま成り立つ。異なるのは絶対速度と細部だけで、
\textbf{並列化すべき対象を分ける境界は微アーキを超えて同じ}である。
Raspberry Pi 5（A76）も同一の枠組みで追記予定である。

% =======================================================================
\section{設計空間 GA：予測 vs 実機実測}
\label{sec:ga}
% ======================================================================
% experiments/ga.tex — 設計空間 GA: 予測 vs 実機実測
% ======================================================================

\subsection{狙い：GA の予測を実測と突き合わせる}

本レポートのここまでは\textbf{実機実測}である。一方、この一連の作業には
別の系譜がある——\texttt{aice-pi-evolution} の\textbf{設計空間 GA}
（\texttt{.aice} → Gemini レビュアー → MAP-Elites → champion）。カーネル設計を
8 軸の遺伝子（scheduler / memory\_alloc / ipc\_primitive / isolation /
arch\_target / power\_mgmt / concurrency\_safety / smp\_model）で表し、
Gemini レビュアー 4 本で採点して世代更新する。

この Round の独自性は、レビュアー rubric を\textbf{本セッションの実機実測に
基づいて}書いたことである。ゆえに GA の champion を「実測で分かっている事実」と
直接突き合わせられる——\textbf{設計空間 GA の予測 vs 実機実測}。

\subsection{結果：worker\_pool が champion（実測が裏付ける）}

seed=4 / gen=10、14 個体。MAP-Elites のセル軸は smp\_model。

\begin{table}[htbp]
\centering
\caption{GA champion（MAP-Elites セル別ベスト）}
\small
\begin{tabular}{@{}lccl@{}}
\toprule
smp\_model & best score & n & 内訳 \\
\midrule
\textbf{worker\_pool} & \textbf{0.622}（champion） & 9 & conc=lock\_free \\
full\_smp & 0.397 & 5 & conc=lock\_free \\
single\_core & (出現せず) & 0 & --- \\
\bottomrule
\end{tabular}
\end{table}

\noindent
\textbf{worker\_pool（0.622）が full\_smp（0.397）に勝った。}これは Pi 5 の
GA Round（0.507 vs 0.371）に続く再現であり、かつ Pi 3 では\textbf{実機実測が
これを裏付けた}——worker\_pool を実装して dining \spd{3.99}・nqueens
\spd{3.66} を実証している。GA の見出し（「worker\_pool を磨け、full\_smp に
走るな」）は、実測の前に、そして独立に、正しい方向を示した。

\subsection{一致とアーティファクト}

champion の 8 軸を実測と照合すると、明確な線が引ける。

\begin{table}[htbp]
\centering
\caption{GA 予測 vs 実機実測}
\small
\begin{tabular}{@{}llll@{}}
\toprule
軸 & GA champion & 実機実測 & 判定 \\
\midrule
smp\_model & worker\_pool & \spd{3.99} 実証 & $\checkmark$ 一致 \\
concurrency\_safety & lock\_free (12/14) & volatile+dsb で成立 & $\checkmark$ 一致 \\
arch\_target & arm\_rpi3 (12/14) & BCM2837 A53 & $\checkmark$ 一致 \\
isolation & mmu\_per\_process (14/14) & none/user\_kernel\_split が安全 & $\times$ アーティファクト \\
ipc\_primitive & capability/tuple\_space & sem\_mailbox（AIPL 保護） & $\times$ アーティファクト \\
\bottomrule
\end{tabular}
\end{table}

\noindent
\textbf{主軸（smp\_model・concurrency\_safety・arch）は実測と一致}する一方、
\textbf{副軸（isolation・ipc）はアーティファクト}である。全 14 個体が
\texttt{isolation=mmu\_per\_process} に固着したのはシード増殖であり、rubric が
減点しても消えなかった。\texttt{ipc=capability/tuple\_space} は実測では
AIPL/JIT/メッシュを全滅させる選択だが champion に残った。これは Pi 5 の
kernel evolution でも見た\textbf{「小規模 GA は方向性まで信頼、全遺伝子詳細は
アーティファクト混入」}パターンの再現である。

\subsection{最大の発見：支配変数が schema に無い}

本セッションの実測で最も効いた変数は\textbf{D-cache}（単一コア \spd{8.8}、
メモリ律速のスケール可否を決める、第 \ref{sec:dcache} 節）だった。しかし
8 軸の遺伝子 schema に\textbf{キャッシュ構成の軸は存在しない}。D-cache は
memory\_alloc / power\_mgmt を通じて間接的にしか表現されず、\textbf{GA は
経験的に最重要と判明した変数を構造的に探索できない}。

したがって具体的な提言が出る：schema に \texttt{cache\_policy} 軸
（\texttt{dcache\_off} / \texttt{dcache\_on\_bounce\_usb} /
\texttt{dcache\_on\_no\_dma}）を追加すべきである。そうすれば GA が D-cache の
トレードオフ（\spd{8.8} vs USB DMA の壁）を直接探索できる。

\subsection{結論：GA は方向を、実測は真実を語る}

\begin{itemize}[leftmargin=1.6em]
  \item \textbf{価値}：GA は「worker\_pool を磨け」という方向を、実測の前に
        独立に正しく示した。Pi 3 / Pi 5 で再現し、Pi 3 では実測が裏付けた。
  \item \textbf{限界 1}：全遺伝子の詳細（isolation・ipc）はアーティファクト
        混入。主軸のみ信頼でき、副軸は人手で実測に合わせて上書きすべき。
  \item \textbf{限界 2}：経験的な支配変数（D-cache）が schema に無ければ、
        GA はそれを見つけられない。実機実測が GA を補完する。
\end{itemize}


% =======================================================================
\section{結論}
\label{sec:conclusion}

\begin{enumerate}[leftmargin=1.6em]
  \item 遊休していた 3 コアを、既存 OS を一切変更せずに計算へ動員できた。
        計算律速の処理で最大 \spd{3.99}（理想値 \spd{4.00}）を実測した。

  \item \textbf{何でも速くなるわけではない}。メモリ書込律速の処理は
        \spd{0.95}--\spd{1.35} に留まる。またプールの往復コスト
        7--11\,\textmu s を下回る小さな仕事は、並列化すると\textbf{遅くなる}
        （実測 \spd{0.38}）。この 2 つの境界が適用対象を決める。

  \item 上記の測定に基づき、\textbf{ウィンドウシステムは並列化しないという
        判断を下した}。描画はメモリ律速であり、かつ 1 回の描画が
        往復コストを下回るため、速くならないばかりか遅くなるからである。
        速くならない変更を入れないことも、工学上の成果である。

  \item AIPL アクターについては、判定基準（1 メッセージあたり 70\,\textmu s
        以上の純計算を持ち、\texttt{send}/\texttt{print}/\texttt{wait} を
        含まない）を満たす唯一の対象に適用した。

  \item 測定基盤そのものにも 3 つの落とし穴があった。タイマ割り込みに依存した
        時計が常に 0 を返していたこと（第\ref{sec:clock}項）、CPU クロックが
        一定でないこと（第\ref{sec:ratio}項）、そして 1 コア基準と並列で
        別関数を測って \spd{4.42} という不可能な値を出していたこと
        （第\ref{sec:rpi3-measbug}項）である。いずれも、比を指標とし、
        フリーランニングカウンタを直読し、両辺で同一関数を測ることで解決した。
        \textbf{安定して再現する誤りは、ノイズより見つけにくい}。

  \item \textbf{支配変数はキャッシュ構成であった}（第 \ref{sec:dcache} 節）。
        D-cache を有効化すると、メモリ律速の処理は作業セットがキャッシュに
        収まる限り理想的にスケールし（\spd{4.01}）、超えると帯域律速へ回帰する
        （\spd{1.40}）。「メモリ律速はスケールしない」は D-cache OFF に固有の
        現象だった。ただし実描画は（キャッシュ不可の FB へキャッシュを超える量を
        書くため）依然スケールせず、ウィンドウシステムの結論は精緻化されつつも
        覆らない。D-cache 有効化には set/way invalidate・SMPEN・DMA コヒーレンシが
        要り、USB がスタック上のバッファを DMA する構造が最大の障害である——
        これは本番投入に残された次段階の課題である。
\end{enumerate}

\end{document}
